% =====================  SECTION V  =====================
\section{Discussion}\label{sec:discussion}

\subsection{Interpretation of Findings}\label{subsec:interpretation}
The results support a nuanced reading of the system. The integration objective is met: a
single user's raw inputs flow, without manual intervention, into an explained,
risk-aligned, and independently verifiable recommendation, and the safety properties hold
by construction rather than by training. What the evaluation also shows is that the
headline classification accuracy is the wrong summary of how the system behaves. The
deterministic guardrail escalates 60.6\% of users to High risk on cash-flow grounds and is
the binding signal for the majority, so the learned model is decisive for only a minority
of users; its 91.0\% accuracy is a real property of one component, measured on held-out
data, but says little about end-to-end decision quality. The comparative and statistical
results reinforce this modest reading: the MLP is statistically indistinguishable from
logistic regression, which in fact edges it on both area-under-curve metrics, so a neural
model is not shown to be necessary. These are not incidental caveats---they are the central
empirical content of the evaluation, and reporting them is part of the contribution.

\subsection{Engineering versus Research Contribution}\label{subsec:contribution-framing}
Separating the two kinds of contribution is what keeps the claims honest. The engineering
contribution is real and is the substance of the work: an integrated pipeline that unifies
feature engineering, unsupervised proxy labeling, a leakage-aware classifier, a
conservative guardrail, most-cautious reconciliation, model-adaptive explanation,
risk-aligned allocation, deterministic simulation, and a hashed, verifiable audit log in
one inspectable system with an explicit feature contract and a safety layer that cannot
lower risk. That combination is exactly what the reviewed literature treats as missing. The
research contribution is stated more modestly: the individual techniques are standard, and
the results do not establish that any is novel or that the learned model beats a simple
baseline. What the work contributes on the research side is methodological---a
demonstration that an integrated financial-recommendation pipeline can be held to an honest
evidential standard in which leakage is measured, guardrail dominance is quantified, and a
model-selection result that fails significance is reported rather than dropped. Ordinary
component integration is claimed nowhere as novel, and no claim is made of validation
against real financial outcomes.

\subsection{Limitations}\label{subsec:limitations}
The critical review surfaced seven limitations, stated here plainly with their severity in
Table~\ref{tab:limitations}. Three are of high severity. Guardrail dominance means the
accuracy figure does not describe end-to-end behavior. Arithmetic label leakage means that,
although the three ratios are withheld from the classifier at the column level, surplus and
expense ratio are deterministic functions of the raw inputs and are therefore recoverable,
making the classification task easier than one with independent labels; the strong
surplus--expense-ratio correlation in Fig.~\ref{fig:corr} is the empirical footprint of
this dependency. Most fundamentally, the labels are proxies produced by clustering, not
recorded default or loss outcomes, so no real-world predictive validity is established. The
medium-severity items are the statistically unjustified selection of the MLP over logistic
regression, the corrected backtest indexing error (which even corrected rests on fixed
synthetic returns), and entity leakage from a row-level split over data with only 944
unique users among 3,000 rows. Finally, at low-to-medium severity, part of the source tree
is a not-yet-extracted presentation layer, the ledger is an in-memory simulation rather
than a deployed chain, and the explainer is occlusion-based rather than SHAP or a
partial-derivative method.

\begin{table}[t]
\centering
\caption{Summary of limitations and their effect on claims}
\label{tab:limitations}
\setlength{\tabcolsep}{4pt}
\begin{tabularx}{\linewidth}{@{}p{3.2cm} >{\centering\arraybackslash}p{1.4cm} L@{}}
\toprule
Limitation & Severity & Effect on claims \\
\midrule
Guardrail dominance & High & Accuracy does not describe end-to-end decision quality. \\
Arithmetic label leakage & High & Classification easier than with independent labels. \\
Proxy labels, no ground truth & High & No real-world predictive validity established. \\
Non-significant model selection & Medium & MLP not statistically better than logistic regression. \\
Backtest indexing error & Medium & Original ROIs inflated; corrected values used. \\
Entity leakage in split & Medium & Generalization may be overstated. \\
Fa\c{c}ade layer / simulated ledger & Low--Med & Prototype-stage engineering, not production. \\
\bottomrule
\end{tabularx}
\end{table}

\subsection{Threats to Validity}\label{subsec:threats}
The limitations map onto the standard categories of validity as follows, so that a reader
can locate each concern precisely.

\begin{itemize}
\item[\textbf{Internal.}] The arithmetic recoverability of two clustering features from the
  classifier inputs, and the row-level split over repeated identifiers, both create paths
  by which apparent performance can be inflated relative to a genuinely independent
  prediction task.
\item[\textbf{Construct.}] The target is a clustering-derived proxy for risk, not a
  recorded financial outcome; the twelve-month ROI is computed from fixed constants rather
  than market data. Both measure the intended construct only indirectly.
\item[\textbf{External.}] The data are synthetic and single-source, and the returns are
  hardcoded, so results should not be extrapolated to real users or markets without
  revalidation on recorded data.
\item[\textbf{Conclusion.}] The small High-risk support (62 holdout cases) widens the
  interval on its per-class F1, and the McNemar tests are underpowered to detect small
  differences; the non-significant results are therefore reported as an absence of
  demonstrated superiority rather than proof of equivalence.
\end{itemize}

\subsection{Research Questions and Objective Fulfillment}\label{subsec:rq-fulfillment}
Table~\ref{tab:rqmap} closes the loop between the six research questions of
Section~\ref{subsec:rqs} and the evidence gathered in Section~\ref{sec:results}, assigning
each an explicit status of Supported, Partially supported, or Not demonstrated together
with the method and result that justify it. The verdicts are deliberately mixed. Three
questions are fully supported: the clustering yields coherent, monotonically ordered proxy
bands (RQ2), the reconciliation enforces an over-warn failure mode by construction (RQ4),
and the allocation-and-audit path produces tamper-evident, re-verifiable records (RQ6). The
remaining three are only partially supported, and are labelled so on purpose---risk can be
assessed from the withheld-ratio inputs, but column-level separation does not remove the
residual arithmetic leakage (RQ1); the classifier reproduces the labels with high accuracy,
yet a neural model is not shown to be warranted over logistic regression (RQ3); and the
explanation mechanism returns per-decision attributions without retraining, but only its
MLP branch is exercised in deployment (RQ5).

\begin{table*}[t]
\centering
\caption{Mapping of research questions to methods, evidence, and status}
\label{tab:rqmap}
\footnotesize
\setlength{\tabcolsep}{5pt}
\renewcommand{\arraystretch}{1.15}
\begin{tabularx}{\textwidth}{@{}>{\centering\arraybackslash}p{0.8cm} L L L >{\raggedright\arraybackslash}p{1.9cm}@{}}
\toprule
RQ & Research question / objective & Method / experiment & Evidence / result & Status \\
\midrule
RQ1 & Assess risk from raw, deployable inputs while withholding the label-defining ratios,
so leakage is explicit. &
Feature contract partitioning $F_{\text{cls}}$ and $F_{\text{clu}}$
(Eq.~\ref{eq:contract}) with a validation abort; residual-leakage analysis
(Sec.~\ref{subsec:contract},~\ref{subsec:limitations}). &
Classifier uses five raw inputs only; column separation is enforced, but surplus and
expense ratio stay arithmetically recoverable, so leakage is measured and disclosed rather
than removed. &
Partially supported \\
\addlinespace
RQ2 & Do interpretable-ratio clusters form coherent, monotonically ordered Low/Medium/High
proxy bands? &
K-means ($k = 3$) on standardized ratios; silhouette, Calinski--Harabasz,
Davies--Bouldin, and ARI stability; cluster profiling (Table~\ref{tab:cluster}). &
Silhouette 0.370 and ARI $0.868 \pm 0.064$; mean surplus falls and mean expense ratio
rises monotonically from Low to High (Table~\ref{tab:cluster}). &
Supported \\
\addlinespace
RQ3 & Can a classifier reproduce the proxy labels from raw inputs with useful accuracy, and
is a neural model warranted over simpler baselines? &
MLP versus logistic regression, random forest, and decision tree; McNemar tests; five-fold
CV; bootstrap CIs (Tables~\ref{tab:perclass}--\ref{tab:stats}). &
91.0\% accuracy and 0.881 macro-F1, but the MLP is not significantly better than logistic
regression ($p = 0.100$) or random forest ($p = 0.176$), which even edge it on AUC. &
Partially supported \\
\addlinespace
RQ4 & Reconcile the learned, unsupervised, and rule-based signals so the failure mode is to
over-warn. &
Escalate-only guardrail (Eq.~\ref{eq:guardrail}) and most-cautious rule
$y^{*} = \max(m, g, k)$ (Eq.~\ref{eq:reconcile}); guardrail-dominance measurement
(Sec.~\ref{subsec:guardrail-dom}). &
Reconciliation cannot lower risk by construction; the guardrail is $\geq$ the cluster
signal for 100\% of users and escalates 60.6\% to High risk. &
Supported \\
\addlinespace
RQ5 & Provide per-decision, input-level attributions across model types without retraining.
&
Model-adaptive occlusion sensitivity (Eq.~\ref{eq:occlusion}): mean-imputation for the MLP,
signed coefficients for linear models, impurity for trees. &
Yields per-input influence scores for the deployed MLP without retraining; the other
branches exist by construction, but only the MLP branch is exercised. &
Partially supported \\
\addlinespace
RQ6 & Map the reconciled band to a risk-aligned crypto allocation and log a tamper-evident,
verifiable record. &
Deterministic band$\rightarrow$(BTC, ETH, USDC) map (Eq.~\ref{eq:alloc}); SHA-256 hashing
of a canonical payload with re-hash verification (Eq.~\ref{eq:hash}). &
Allocation is deterministic and risk-monotone; any change to a logged payload changes its
digest and is detected---though against an in-memory simulated ledger, not a live chain. &
Supported \\
\bottomrule
\end{tabularx}
\par\vspace{3pt}
{\footnotesize\itshape\raggedright The three `supported' verdicts rest on properties that
hold by construction (RQ4, RQ6) or on direct metric evidence (RQ2); the three `partially
supported' verdicts mark exactly where an honest reading stops short---residual leakage, an
unjustified neural model, and a single exercised explanation branch. No verdict claims
validation against real financial outcomes, which remains outside the scope of this
proxy-labelled study.\par}
\end{table*}
